% Chapter 2, section 6: Retries and the dead-letter queue
% Write in plain LaTeX. Code of a guide step: \guide{<step>}; commands: \begin{shell}...\end{shell};
% a screenshot: \screenshot{file}{caption}; a gap in the draft: \placeholder{...}.
% See book/latex/queen.sty for the full list.

\section{Retries and the dead-letter queue}

Every \texttt{failed} ack spends one retry. The message comes back after the queue's
\texttt{retryDelay}, and its partition waits with it, for that consumer group. After
\texttt{retryLimit} failures, the message goes to the dead-letter queue and the group moves on. You
set both on the queue (see Create and configure queues).

\why{retrying is how temporary failures heal by themselves, and the limit is how a message that can
never succeed stops blocking everything behind it in its partition.}

Two cases do not spend a retry: a \texttt{retry} ack, and a lease that runs out. So a message that
crashes its worker's process every time comes back forever and never reaches the dead-letter queue.
If some input can crash your worker, check for it before the work starts and ack such a message with
\texttt{dlq}.

\queenfig{retry-dlq}{The life of a message: four endings, and the replay.}

\subsection{What a dead letter keeps}

A dead letter keeps the payload, the error the worker sent, the number of retries, and the consumer
group that gave up on it. It stays until you replay it or delete it, or until its queue holds too
many: past 100,000 dead letters, 1 GiB or 14 days in one queue, the oldest go first (from broker
2.7.0; each cap is a setting, and \texttt{0} turns it off). Each group has its own: a message that
failed for billing can be fine for analytics.

\why{a failure you cannot see is a failure you cannot fix. The dead-letter queue is where they wait
for a person, or for a fix in the code.}

\subsection{Read it}

\guide{dlq-read}

\texttt{limit()} and \texttt{offset()} page through it, newest first. The client also has
\texttt{from()} and \texttt{to()}, but the 2.7.0 broker ignores them.

\subsection{Replay it}

Once the cause is fixed, put the messages back:

\guide{dlq-replay}

A replay moves the dead letter back into its partition, at the end, as a new message with the
\texttt{transactionId} \texttt{dlq:<id>}, and removes it from the dead-letter queue in the same
step. Doing it twice cannot create a second copy. Every group that reads that partition sees the
replayed message, not only the group that gave up on it. Over HTTP the same thing is
\texttt{POST /api/v1/dlq/:id/replay}.

\why{the fix usually comes after the failure: a bug deployed, a provider back up. Replay lets the
same messages go through again without anyone pushing them by hand.}

\subsection{Skip it, or turn it off}

An ack with \texttt{dlq} files a message straight away, without spending its retries. Setting
\texttt{deadLetterQueue} and \texttt{dlqAfterMaxRetries} to false on the queue drops a message
that runs out of retries instead of keeping it.

\why{some messages are worth nothing once they fail (a typing indicator, a stale price), and
keeping them would only fill a list nobody reads.}

On the site, with the code in all six SDK languages: \href{https://queenmq.com/writing-applications/retries/}{queenmq.com/writing-applications/retries}. More there: \href{https://queenmq.com/writing-applications/pop/}{consuming}.
